% Vietnamese translation for OI Public Library
% Source-File: IOI中国国家候选队论文/2004/栗师.ppt
\documentclass[11pt]{article}
\usepackage{oipl-vn}

\title{Chuyển hóa mục tiêu trong giải bài\\{\large Ghi chú theo slide}}
\author{Li Shi}
\date{IOI 2004}

\begin{document}
\maketitle

\origtitle{Transforming goals in problem solving}
\sourcepath{IOI 2004 / Li Shi.ppt}

\section{Ý tưởng chung}

Bài trình bày nói về việc thay mục tiêu khó bằng một mục tiêu gần hơn, dễ
phân tích hơn, rồi quay lại mục tiêu ban đầu. Cách làm này không phải đoán
mò: mục tiêu mới phải giữ được thông tin cốt lõi, hoặc đủ mạnh để suy ra lời
giải cho bài gốc.

\section{Mốc từ slide}

Phần văn bản đọc được nhắc các vector bước đi
\[
  (x,y)\in\{(0,1),(0,-1),(1,0),(-1,0)\},
\]
các trạng thái \(S_{i,j}\), \(T_{i,j}\), và tổ hợp tuyến tính
\[
  k_1P_i+k_2P_j=(0,y).
\]
Đây là các ký hiệu trong ví dụ về một quân cờ trên lưới vô hạn: thay vì hỏi
ngay liệu có đi tới mọi ô hay không, ta xét nhóm sinh bởi các vector bước đi,
đưa bài toán về điều kiện số học trên các tổ hợp nguyên.

\section{Cách dùng trong thiết kế thuật toán}

Khi một bài có quá nhiều ràng buộc, ta có thể:
\begin{enumerate}
  \item bỏ bớt một ràng buộc để tìm cấu trúc chính;
  \item mở rộng hoặc thu hẹp miền cần xét;
  \item chứng minh mục tiêu mới tương đương, hoặc đủ để phục hồi mục tiêu cũ.
\end{enumerate}
Điểm then chốt là luôn ghi rõ cầu nối giữa hai mục tiêu. Nếu cầu nối yếu,
thuật toán có thể giải rất tốt một bài khác nhưng không giải bài đang xét.

\end{document}
